\section{A squared-modulus quotient: a 66-operation equation}
\label{sec:squared-modulus}

The constructions with fewest variables need large arithmetic circuits.
There is a separate improvement at the short-circuit end of the comparison:
the inherited fifteen-positive-variable equation can be evaluated in
\textbf{66 binary arithmetic operations}.  Its exact degree is 157.  This
saves one operation against the repository's 67-operation, degree-155
baseline.  The same idea lowers the costs of the degree-80 and degree-68
alternatives from 73 and 75 to 72 and 74 operations, respectively.
The saving comes from a stronger auxiliary congruence, followed by a change
of the quantified quotient.  It is not an identity between the old and new
polynomials at the same quotient coordinate.

\subsection{The explicit polynomial and its fixed coefficients}

Use the entire coefficient tuple of package 36 at the source commit
\texttt{124c07c0b7}: the ports are
\(\beta=B-1,K_0,M_C,M_F\), with \(B=2^{1012115}\).
Here \(K_0\) is the tuple's \texttt{K} and \(M_F\) is its already shifted
\texttt{MF\_source}.  The four ports are fixed integers, not additional
unknowns.  The coefficient construction and its machine semantics are
inherited; the argument below introduces no rebinding or change of masks.

All fifteen unknowns are strictly positive integers:
\[
 (J,F,\alpha,z,f,h,i,T,s,w,\tau,\eta,\zeta,y,\gamma)\in\Pos^{15}.
\]
Define the following polynomial abbreviations, which introduce no new
unknowns:
\begin{align*}
 q&=\beta J+1,                 & C&=q-F-\alpha,\\
 X&=wq,                       & Y&=sq^3, & E&=XY,\\
 k&=\eta+\zeta,               & c&=kY+\eta, & a&=E+Y,\\
 H&=4a+3,                     & \Delta&=a^2+H, & A&=a+2,\\
 D&=X+ac+\gamma H,            & G&=(q-1)(q-F)+\alpha,\\
 R&=G(q^2-1)+(M_C+qM_F)J,     & S&=i\Delta c^2,\\
 V&=f^2(cT-R)-c.&&&
\end{align*}
In particular, \(\Delta=A^2-1=(a+1)(a+3)>0\) throughout the positive
domain, before any equation is imposed.  Set
\begin{align*}
 N_0&=\tau^2-(XY^2)(XY^2+1)k^2,\\
 N_1&=D^2-\Delta c^2,\\
 N_a&=S^2(V^2-y^2)+y^2,\\
 N_k&=k-hE-R,\\
 N_t&=(K_0+w)C+(q-F)-z(q-1),\\
 N_{\mathrm{str}}&=\Delta f^2-S^2.
\end{align*}
The proposed equation is the ordinary integer polynomial equation
\begin{equation}
 \boxed{\quad
 P_{66}:=N_0N_1N_aN_kN_tN_{\mathrm{str}}-\Delta=0.
 \quad}
 \label{eq:p66}
\end{equation}
Every exponent in this definition is fixed.  The divisions and Pell
polynomials used below occur only in the proof and witness construction.

\subsection{Why a square modulus permits a cheaper quotient}

Write \(\chi_t(n),\psi_t(n)\) for the integer Pell polynomials determined by
\[
 \chi_t(0)=1,\qquad \psi_t(0)=0,
\]
\[
 \chi_t(n+1)=t\chi_t(n)+(t^2-1)\psi_t(n),\qquad
 \psi_t(n+1)=\chi_t(n)+t\psi_t(n).
\]
Thus \(\chi_t(n)^2-(t^2-1)\psi_t(n)^2=1\).

\begin{lemma}[Odd Pell quotient identities]
\label{lem:odd-quotient}
For each \(r\in\N\) there is a polynomial \(Q_r\in\Z[Z]\) such that
\begin{align}
 \chi_t(2r+1)&=tQ_r(t^2),\label{eq:odd-quotient-identity}\\
 Q_r(0)&=(-1)^r(2r+1),\label{eq:odd-quotient-constant}\\
 Q_r(1-A^2)&=(-1)^r\psi_A(2r+1).
 \label{eq:odd-quotient-conjugate}
\end{align}
\end{lemma}

\begin{proof}
The odd-index subsequence gives
\[
 Q_0(Z)=1,\qquad Q_1(Z)=4Z-3,\qquad
 Q_{r+1}(Z)=(4Z-2)Q_r(Z)-Q_{r-1}(Z).
\]
This proves \eqref{eq:odd-quotient-identity} and
\eqref{eq:odd-quotient-constant} by induction.  The sequence
\(\psi_A(2r+1)\) starts with \(1,4A^2-1\) and has recurrence coefficient
\(4A^2-2\).  Substituting \(Z=1-A^2\) changes the coefficient of the
displayed recurrence to its negative, giving
\eqref{eq:odd-quotient-conjugate} by the same induction.  These are
identities over the integers and therefore also hold modulo any positive
integer.
\end{proof}

\begin{lemma}[Squared-modulus reconstruction]
\label{lem:squared-modulus-reconstruction}
Suppose \(A>1\), \(R>0\), \(R\equiv3\pmod4\),
\(c=\psi_A(R)>0\), and \(i,f>0\) satisfy
\[
 \Delta=A^2-1,\qquad f^2=1+\Delta i^2c^4.
\]
Put \(S=\Delta ic^2\).  There are positive integers \(V,y,T\) satisfying
\[
 S^2(V^2-y^2)+y^2=1,\qquad V=f^2(cT-R)-c.
\]
\end{lemma}

\begin{proof}
Take
\[
 V=\frac{\chi_S(R)}{S},\qquad y=\psi_S(R).
\]
The quotient is integral by \eqref{eq:odd-quotient-identity}; both values
are positive because \(S>1\) and \(R>0\).  The Pell identity gives the
required norm equation.  Also
\[
 S^2=\Delta(f^2-1)\equiv1-A^2\pmod{f^2}.
\]
Since \((R-1)/2\) is odd, Lemma~\ref{lem:odd-quotient} gives the two
congruences
\begin{equation}
 V\equiv-c\pmod{f^2},\qquad V\equiv-R\pmod c.
 \label{eq:squared-modulus-congruences}
\end{equation}
For the second, use \(c\mid S\) and
\eqref{eq:odd-quotient-constant}.  The strong equation gives
\(f^2\equiv1\pmod c\), hence \(\gcd(c,f)=1\).
The numerator \(V+c+Rf^2\) is therefore divisible by both \(f^2\) and
\(c\), and thus by \(cf^2\).  Consequently
\[
 T=\frac{V+c+Rf^2}{cf^2}
\]
is a positive integer.  Rearranging proves the stated quotient identity.
\end{proof}

The earlier auxiliary construction used congruence modulo \(f\).  Retaining
the full modulus \(f^2\) is what makes the old quotient divisible by
\(f\), after choosing a suitable auxiliary Pell ordinate.

\subsection{Soundness, completeness, and the native congruence}

\begin{theorem}[Existence-equivalence of the short equation]
\label{thm:p66-consistency}
With the authentic package-36 coefficient tuple, the inherited ordinary
mathematical compiler and K215 calibration imply
\[
 \mathrm{ZFC}\vdash
 \left[
 \neg\exists\mathbf{x}\in\Pos^{15}\;P_{66}(\mathbf{x})=0
 \right]
 \;\longleftrightarrow\;
 \operatorname{Con}(\mathrm{ZFC}).
\]
The additional reduction from the inherited explicit polynomial to
\eqref{eq:p66} requires no consistency assumption.
\end{theorem}

\begin{proof}
Let \(P_{67}\) denote the inherited polynomial on the same fifteen slots.
Its only different abbreviation is
\[
 V_{\mathrm{old}}=c(T_{\mathrm{old}}f-1)-Rf^2.
\]
The exact all-value identity
\begin{equation}
 P_{66}(\ldots,f,T,\ldots)
   =P_{67}(\ldots,f,fT,\ldots)
 \label{eq:p66-pullback}
\end{equation}
follows by substitution.  Thus every positive new zero gives a positive old
zero, since \(fT>0\).  Inherited soundness supplies a finite ZFC refutation.

Conversely, a finite refutation supplies a positive old zero.  Its inherited
Pell and compiler recovery gives
\[
 c=\psi_A(R),\qquad f^2=1+\Delta i^2c^4,
 \qquad q=B^N,\qquad J=1+B+\cdots+B^{N-1}
\]
for some positive integer \(N\), together with the native bit condition
\[
 (C-1)\mathbin{\mathrm{AND}}(M_CJ+1)=0.
\]
The tuple has \(4\mid B\) and \(M_C\equiv2\pmod4\).  Hence
\(q\equiv0\pmod4\), \(J\equiv1\pmod4\), and
\(M_CJ+1\equiv3\pmod4\).  The two low bits of the nonnegative integer
\(C-1\) must vanish, so \(C\equiv1\pmod4\).  Using the exact identity
\(G=q(q-F)-C\), the definition of \(R\) yields
\[
 R\equiv C+2\equiv3\pmod4.
\]
This phase follows from the authentic native decoding; it is not a
restriction silently imposed on arbitrary coefficient assignments.

Apply Lemma~\ref{lem:squared-modulus-reconstruction} at these same values
of \(A,R,c,i,f\).  Replace only the supplied auxiliary ordinate \(y\) and
quotient \(T\) by the lemma's positive integers.  The auxiliary norm factor
is again one.  The other factors are unchanged; in particular the inherited
unit classification gives \(N_0=N_1=N_k=N_t=1\) and
\(N_{\mathrm{str}}=\Delta\).  Thus \(P_{66}=0\).

We have established equivalence of positive solvability, which composes
with the inherited refutation calibration.  Negating gives the displayed
consistency equivalence.  All new steps are finite integer algebra,
induction on polynomial recurrences, and elementary divisibility, and can
be carried out in ZFC.
\end{proof}

The reverse map may change both auxiliary coordinates.  In particular,
the theorem does not assert that every old quotient is divisible by \(f\),
or that the old and new equations have identical zero sets at unchanged
coordinates.  The local substitution proof also does not supply an
end-to-end Lean or Rocq certificate for the inherited native-window
compiler.

\subsection{Exact arithmetic size, degree, and companion equations}

The old evaluator already computes \(f^2\).  Its five private quotient
operations are replaced by the four operations
\[
 cT,\qquad cT-R,\qquad f^2(cT-R),\qquad f^2(cT-R)-c.
\]
Every other gate is retained.  The complete new evaluator has
\(38\) multiplications, \(15\) additions, and \(13\) subtractions.
Multiplication by a fixed numeral is charged; reuse is free.  All gates,
all fifteen witnesses, and all four coefficient ports are live.

The exact factor degrees in the order
\(N_0,N_1,N_a,N_k,N_t,N_{\mathrm{str}}\) are
\[
 22,\quad18,\quad62,\quad7,\quad2,\quad46.
\]
Put \(Q=\beta J\), \(k_0=\eta+\zeta\), and
\(C_1=Q-F-\alpha\).  The leading homogeneous polynomial of \(P_{66}\) is
\[
 -8Q^{91}h\gamma i^4k_0^{13}w^{13}s^{26}T^2f^4
       (wC_1-zQ).
\]
In particular, the coefficient of
\[
 J^{92}h\gamma i^4\eta^{13}w^{13}s^{26}T^2f^4z
\]
is \(8\beta^{92}\ne0\).  The exact degree is therefore 157, and the
final subtraction of \(\Delta\), of degree 12, cannot affect it.

For lower-degree alternatives, retain the new \(V\) but put
\[
 \widetilde N_a=\Delta(f^2-1)(V^2-y^2)+y^2,
 \qquad N_s=f^2-\Delta(ic^2)^2.
\]
The following three companion equations inherit the same positive
solvability equivalence:
\begin{align*}
 P_{68}&=N_0N_1\widetilde N_aN_kN_tN_s-1=0,\\
 P_{72}&=(N_0N_1-1)^2+(\widetilde N_aN_kN_t-1)^2+(N_s-1)^2=0,\\
 P_{74}&=(N_0-1)^2+(\widetilde N_a-1)^2+(N_s-1)^2
              +(N_1N_kN_t-1)^2=0.
\end{align*}
Indeed each is the corresponding inherited parent equation under the same
positive substitution \(T_{\mathrm{old}}=fT\); the established parent
equivalences and the reconstruction proof above apply.  Their fully paid
counts and exact degrees are
\[
\begin{array}{c|r|r|r|r|r}
 &\text{multiplications}&\text{additions}&\text{subtractions}
 &\text{total}&\text{degree}\\ \hline
 P_{66}&38&15&13&66&157\\
 P_{68}&39&15&14&68&113\\
 P_{72}&40&17&15&72&80\\
 P_{74}&40&18&16&74&68
\end{array}
\]
For \(P_{72}\), the first squared residual uniquely reaches degree 80;
its leading form is
\(64\gamma^2w^8s^{14}(\eta+\zeta)^6(\beta J)^{50}\).
For \(P_{74}\), the strong residual uniquely reaches degree 68 after
squaring; its leading form is
\(i^4w^4s^{12}(\eta+\zeta)^8(\beta J)^{40}\).
For \(P_{68}\), a degree-113 monomial has coefficient
\(8\beta^{64}\).  These nonzero coefficients certify the claimed exact
degrees uniformly for the prescribed tuple.

The supplied standard-library auditor checks the literal schedules, source
hashes, exact polynomial identities at the changed cut, leading homogeneous
forms, and full liveness.  Normal and optimized Python executions agree.
Its finite congruence examples and modular circuit diagnostics support the
implementation; the universal reconstruction lemma, rather than those
examples, supplies the new mathematical proof.  The result improves stated
repository tradeoffs and makes no claim of global minimality.
